\par
\textsc{Gauss} hat also auch die Mehrdeutigkeit eines Integrals einer \textls[100]{eindeutigen
Funktion}, aufgefasst als Funktion der oberen Grenze, bei der Integration
im komplexen Gebiete gekannt. Dagegen scheint er die Integration
einer mehrdeutigen Funktion im komplexen Gebiete, namentlich die Theorie
der Verzweigungspunkte nicht gekannt und auch für die Erscheinung, dass
ein Integral zwischen bestimmten Grenzen von der Form \(\int_\alpha^\beta\varphi(z,x)\,dz\) eine unendlich
vieldeutige Funktion des Parameters \(x\) sein kann, keine Erklärung
besessen zu haben. Diese Annahme würde es begreiflich erscheinen lassen,
dass er die Umkehrungsfunktion des elliptischen Integrals erster Gattung
nicht \textls[100]{direkt} studiert, ferner, dass er für die Perioden, also für das agM.
die Darstellung durch das vollständige elliptische Integral, und allgemeiner\par
\clearpage\SourceLeaf{238}\ArithmeticHead{153}{R}{VIId. FUNKTIONENTHEORETISCHES.}
\noindent für \(F(\alpha,\beta,\gamma,x)\) die \textsc{Eulersche} Darstellung durch das Integral
\[
\int_0^1 z^{\beta-1}(1-z)^{\gamma-\beta-1}(1-xz)^{-\alpha}\,dz
\]
nicht benutzt hat, während er — wie oben erwähnt wurde — z. B. von dem
sogenannten \textsc{Eulerschen} Integrale zweiter Gattung (12), das eine \textls[100]{eindeutige}
Funktion des Parameters ist, Gebrauch macht.
\par
Aber noch ein anderes geht aus der Bemerkung des Art. 3 hervor! \textsc{Gauss}
braucht, um die durch \(F(\alpha,\beta,\gamma,x)\) für \(|x|<1\) definierte Funktion fortzusetzen,
»ein höheres Prinzip«, das ihm durch die Differentialgleichung geliefert wird.
Vermöge dieser erfolgt dann die Fortsetzung nach dem Gesetze der \textls[100]{Kontinuität}.
Aber er scheint nicht gewusst zu haben, dass ein solches höheres
Prinzip überflüssig oder vielmehr von selbst wirksam ist, wenn es sich um
eine Funktion einer komplexen Variabeln handelt, d. h. es scheint ihm damals
das Prinzip der \textls[100]{analytischen} Fortsetzung gefehlt zu haben. Aus dem
vielleicht unbestimmten, aber immerhin vorhandenen Bewusstsein, dass seine Einsichten
in die Natur der Funktionen einer komplexen Veränderlichen noch Lücken
aufweisen, liesse sich die Zurückhaltung erklären, die \textsc{Gauss} der Öffentlichkeit
gegenüber mit seinen analytischen Arbeiten stets bewahrt hat, und man muss
auch sagen, dass es ihm geradezu unmöglich gewesen wäre, die Gesamtheit
seiner Untersuchungen über die transzendenten Funktionen zu einem harmonischen
Ganzen zu gestalten, wenn er wirklich, um es kurz zu bezeichnen,
in der Theorie der Integrale nicht über den Standpunkt \textsc{Cauchys} hinausgekommen
war.
\par
\setcounter{footnote}{216}
P. \textsc{Günther} berichtet\footnote{P. \textsc{Günther}, Göttinger Nachrichten 1894, S. 104.} über eine ihm von \textsc{Weierstrass} mitgeteilte
Äusserung \textsc{Wilhelm Webers}, wonach \textsc{Gauss} nicht eher an die Veröffentlichung
seiner Untersuchungen über die transzendenten Funktionen gehen wollte, »als
bis er die verschiedenen Methoden, die ihm den Eingang in ihre Theorie vermittelt
hatten, mit einander vollkommen in Einklang gesetzt haben würde«.
Das wäre ihm aber nur möglich gewesen, wenn er bis zu dem funktionentheoretischen
Standpunkte \textsc{Riemanns} vorgedrungen wäre, und das scheint auch
später nicht der Fall gewesen zu sein, \textls[100]{wenigstens was die Theorie der
Integrale mehrdeutiger Funktionen anlangt}\footnote{Auf \textsc{Gauss'} Auffassung des Integrals wird im Abschnitt X noch einmal zurückzukommen sein.}. In bezug auf das\par
\HistoricalFoot{X 2 Abh. 2.}{20}
\clearpage\SourceLeaf{239}\ArithmeticHead{154}{L}{L. SCHLESINGER, ÜBER GAUSS' FUNKTIONENTHEORETISCHE ARBEITEN.}
\noindent \textls[100]{Prinzip der analytischen Fortsetzung} hat \textsc{Gauss} später alle Hilfsmittel
bereit gestellt, die zur Formulierung dieses Prinzips erforderlich sind; wir
kommen auch darauf weiter unten zurück.\par

